\begingroup\small
\begin{thebibliography}{99}
\raggedright
\bibitem{GS} S. B. Giddings and A. Strominger, Axion-induced topology change in quantum gravity and string theory, \emph{Nucl. Phys. B} \textbf{306} (1988) 890--907. \href{https://doi.org/10.1016/0550-3213(88)90446-4}{doi:10.1016/0550-3213(88)90446-4}.
\bibitem{Kaul} R. K. Kaul and S. Sengupta, Torsional instanton effects in quantum gravity, \emph{Phys. Rev. D} \textbf{90} (2014) 124081, \href{https://arxiv.org/abs/1411.4793}{arXiv:1411.4793}.
\bibitem{Kim} H. Kim, Axionic wormholes: more on their classical and quantum aspects, \emph{Nucl. Phys. B} \textbf{527} (1998) 311--341, \href{https://arxiv.org/abs/hep-th/9706015}{hep-th/9706015}.
\bibitem{HTV} T. Hertog, B. Truijen and T. Van Riet, Euclidean axion wormholes have multiple negative modes, \emph{Phys. Rev. Lett.} \textbf{123} (2019) 081302, \href{https://arxiv.org/abs/1811.12690}{arXiv:1811.12690}.
\bibitem{LSS} G. J. Loges, G. Shiu and N. Sudhir, Complex saddles and Euclidean wormholes in the Lorentzian path integral, \emph{JHEP} 08 (2022) 064, \href{https://arxiv.org/abs/2203.01956}{arXiv:2203.01956}.
\bibitem{MM} D. Marolf and B. Missoni, Euclidean wormholes stability analysis revisited (2025), \href{https://arxiv.org/abs/2505.21118}{arXiv:2505.21118}.
\bibitem{TY} S. Takeshita and N. Yoshioka, Axionic wormholes in metric-affine gravity (2026), \href{https://arxiv.org/html/2609.11010v2}{arXiv:2609.11010v2}.
\bibitem{P08} Muacca. Reading the sign of a Nieh--Yan charge on compact spin three-manifolds. v1 (2026). \href{https://doi.org/10.5281/zenodo.24434717}{doi:10.5281/zenodo.24434717}.
\bibitem{BarSpin} C. B\"ar, The Dirac operator on space forms of positive curvature, \emph{J. Math. Soc. Japan} \textbf{48} (1996) 69--83, \href{https://www.jstage.jst.go.jp/article/jmath1948/48/1/48_1_69/_article/-char/en}{J-STAGE}.
\bibitem{BB} C. B\"ar and W. Ballmann, Guide to boundary value problems for Dirac-type operators, \href{https://arxiv.org/abs/1307.3021}{arXiv:1307.3021}.
\bibitem{KY} S. K. Kobayashi and K. Yonekura, The Atiyah--Patodi--Singer index theorem from the axial anomaly, \emph{PTEP} (2021) 073B01, \href{https://doi.org/10.1093/ptep/ptab061}{doi:10.1093/ptep/ptab061}.
\bibitem{LV} M. Lesch and B. Vertman, Regular singular Sturm--Liouville operators and their zeta-determinants, \emph{J. Funct. Anal.} \textbf{261} (2011) 408--450, \href{https://arxiv.org/abs/1005.0368}{arXiv:1005.0368}.
\bibitem{LVsum} M. Lesch and B. Vertman, Regularizing infinite sums of zeta-determinants, \emph{Math. Ann.} \textbf{361} (2015) 835--862, \href{https://arxiv.org/abs/1306.0780}{arXiv:1306.0780}.
\bibitem{DF} X. Dai and D. S. Freed, Eta-invariants and determinant lines, \emph{J. Math. Phys.} \textbf{35} (1994) 5155--5194, \href{https://arxiv.org/abs/hep-th/9405012}{hep-th/9405012}.
\bibitem{WY} E. Witten and K. Yonekura, Anomaly inflow and the eta-invariant, \href{https://arxiv.org/abs/1909.08775}{arXiv:1909.08775}.
\bibitem{KSZ} G. K. Karananas, M. Shaposhnikov and S. Zell, On the gravitational origin of the QCD axion, \emph{Phys. Rev. Lett.} \textbf{135} (2025) 241001, \href{https://arxiv.org/abs/2506.11836}{arXiv:2506.11836}.
\end{thebibliography}
\endgroup
